\chapter{Danksagung}
\selectlanguage{ngerman}
Die vorliegende Dissertation entstand während meiner Tätigkeit als wissenschaftlicher Mitarbeiter
am F. A. Finger-Institut für Baustoffkunde, Professur Werkstoffe des Bauens an der Bauhaus-Universität Weimar im Rahmen mehrerer, von der DFG geförderter Projekte.

An dieser Stelle möchte ich mich bei allen bedanken, die zum Gelingen dieser Arbeit beigetragen haben. Mein herzlicher Dank gilt:

\begin{itemize}
    \item Professor Dr.-Ing. Horst-Michael Ludwig für die Möglichkeit, meine Dissertation am F. A. Finger-Institut für Baustoffkunde an der Bauhaus-Universität Weimar zu schreiben,
    \item Dr. Christiane Rößler für die intensive fachliche Unterstützung und Motivation,
    \item M. Sc. Sophie Unbehau für die Unterstützung bei der Messung an der \acs*{LTDSC} und die gemeinsame Auswertung der Ergebnisse,
    \item Prof. Harald Hilbig und Dr. Marco Decker der TU München für die \acs*{LAICPMS}-Messungen und die Unterstützung bei der Auswertung,
    \item Prof. Dr. Jürgen Neubauer und M.Sc. Franz Becker der FAU Erlangen-Nürnberg für die Anfertigung einiger \acs*{XRD}-, \acs*{XRF}- und \ce{^{1}H}-NMR-Messungen,
    \item M.Sc. Feng Li der Materialforschungs- und Prüfanstalt (MFPA) Weimar für das Zurverfügungstellen von Proben und die Nanoindentationsversuche,
    \item Prof. Dr. Thomas Höche und Dipl.-Ing. Georg Schusser des Fraunhofer-Instituts für Mikrostruktur von Werkstoffen und Systemen (IMWS) für die Laser-Präparation einiger Proben,
    \item den ehemaligen Studierenden Melanie Heinemann, Paul Schilder und Inga Rust, deren studentische Arbeiten meine Dissertation unterstützt haben.
\end{itemize}

Weiterhin danke ich den Angehörigen des F. A. Finger-Instituts für Baustoffkunde für viele hilfreiche Diskussionen, insbesondere Dipl.-Ing. Christian Matthes für die gemeinsamen Untersuchungen an den Rasterelektronenmikroskopen, sowie Lydia Schmiedel und Nadine Brendel für die Politur zahlreicher Proben.

Abschließend danke ich meinen engen Freunden Dr. Michael Böhme, M. Sc. Frowin Ellermann und M. Sc. Sophie Unbehau die mich auf dem Weg meines Promotionsstudiums immer wieder zur Arbeit an diesem Dokument motivieren konnten.
